\section{Power corrections}\label{sect_powercorrections}%
%

A singularity on the integration path in Eq.~\eqref{BorelIntegral} means that the perturbation theory is 
incomplete and the sum of the series is ill-defined. 
%Following ~\cite{Beneke:1998ui,Beneke:2000kc}
It is customary ~\cite{Beneke:1998ui} to estimate the corresponding ambiguity as
\begin{align}
\delta H(w_0)  = -\pi \frac{1}{\beta_0}  e^{-w_0/(\beta_0 a_s)} \res_{w=w_0}\big[B[H](w)\big]\,, 
\end{align}
where $w_0$ is the position of the singularity and $\displaystyle{\res_{w=w_0}\big[B[H](w)\big]}$ is the corresponding residue. Note that  $ e^{-w_0/(\beta_0 a_s)} = (\Lambda^2/\mu^2)^{w_0}$. Following the standard logic~\cite{Beneke:1998ui,Beneke:2000kc} we assume that this ambiguity must be 
canceled by adding a non-perturbative correction of the same order of magnitude. 

%
\subsection{Ioffe time quasi-distributions}\label{sect_qITD}%
%

Considering $\delta H(1)$, we obtain the leading power correction to the qITDs as functions of the ``Ioffe-time''
\begin{align}
 \tau = (p\cdot v) z 
\end{align}
\begin{align}
\label{t4-qITD}
  \mathcal{I}^\parallel(\tau) 
&= I(\tau) +
  \kappa (v^2\! z^2\!\Lambda^2)\! \int_0^1\!d\alpha \, (\alpha + \bar\alpha \ln \bar\alpha)I(\alpha \tau )\,,
\notag\\
  \mathcal{I}^\perp(\tau)
%(z^2v^2, pvz, \mu) 
&= I(\tau) +
  \kappa (v^2\! z^2\!\Lambda^2) \!\int_0^1\!d\alpha \, (\alpha + \bar\alpha \ln \bar\alpha + \alpha\bar\alpha)I(\alpha \tau )\,,
\end{align}
where $\kappa$ is a real number of order one. The renormalon ambiguity \eqref{w=1} corresponds to 
\begin{align}
  |\kappa| = -\pi \frac{1}{\beta_0}\left(- \frac{ e^{5/3}}{4 e^{-2\gamma_E}}\right) = \frac{\pi e^{5/3+2\gamma_E}}{4 \beta_0} \simeq 1.5\,,  
\end{align} 
but this number is only indicative. Alternatively, one can put $\kappa=1$ and think of $\Lambda$ as a certain nonperturbative 
parameter of the order of $\Lambda_{\rm QCD}$ that determines an overall normalization of the power correction and 
cannot be fixed in this approach. Note that also the sign of the correction is not determined.

For the normalized qITDs defined in Eq.~\eqref{n-qITD} we obtain instead
\begin{align}
\label{t4-nqITD}
  \mathbf{I}^\parallel(\tau)
%(z^2v^2, pvz, \mu) 
&= I(\tau) +
  \kappa (v^2\! z^2\!\Lambda^2) \!\int_0^1\!d\alpha \, [\alpha + \bar\alpha \ln \bar\alpha]_+I(\alpha \tau )\,,
\notag\\
  \mathbf{I}^\perp(\tau)
%(z^2v^2, pvz, \mu) 
&= I(\tau) +
  \kappa (v^2\! z^2\!\Lambda^2) \!\int_0^1\!\!d\alpha \, [\alpha + \bar\alpha \ln \bar\alpha + \alpha\bar\alpha]_+I(\alpha \tau )\,,
\end{align} 
where the ``plus'' distribution is defined as usual,
\begin{align}
  [f(\alpha)]_+ = f(\alpha) - \delta(\bar\alpha)\int_0^1 d\beta\, f(\beta)\,. 
\end{align}
The leading power corrections to the qITDs 
$\widehat{\mathbf{I}}^{\parallel(\perp)}$ ~\eqref{hat-qITD} that are normalized to the vacuum correlator,
are the same as for the un-normalized distributions \eqref{t4-qITD}.
The expressions in~\eqref{t4-qITD} and \eqref{t4-nqITD} present the starting point for the following analysis.

In order to visualize the functional dependence of the power correction on the ``Ioffe time'' $\tau$ relative to the leading-twist result, we write 
\begin{align}
    \cI &= I(\tau)\Big\{ 1 + \kappa (v^2 z^2\Lambda^2) \cR_{\cI}(\tau)\Big\},
\notag\\
    \bI &= I(\tau)\Big\{ 1 + \kappa (v^2 z^2\Lambda^2) \bR_{\cI}(\tau)\Big\},
\end{align}
where
\begin{align}
\cR_{\cI}^\parallel(\tau) & = \frac{1}{I(\tau)} \int_0^1\!d\alpha \, 
(\alpha + \bar\alpha \ln \bar\alpha)I(\alpha \tau)\,,
\notag\\  
\cR_{\cI}^\perp(\tau)
& = \frac{1}{I(\tau)} \int_0^1\!d\alpha \, 
(\alpha + \bar\alpha \ln \bar\alpha + \alpha\bar\alpha)I(\alpha \tau)\,.
\end{align}
For the normalized qITDs one obtains 
\begin{align}
  \bR_{\cI}^\parallel(\tau) &= \cR_{\cI}^\parallel(\tau) - \frac14\,, 
\qquad
  \bR_{\cI}^\perp(\tau) &= \cR_{\cI}^\perp(\tau) - \frac5{12}\,, 
\end{align}
and, obviously, 
\begin{eqnarray}
\widehat{\bR}_{\cI}^{\parallel(\perp)}(\tau) = \cR_{\cI}^{\parallel(\perp)}(\tau),
\end{eqnarray}
so that we do not consider them separately.


In general, the qITDs $\cI(\tau)$ and the higher-twist coefficients $\cR_{\cI}(\tau)$ and $\bR_{\cI}(\tau)$ are complex functions, but their imaginary parts appear to be small. The real parts, $\Real\cR_{\cI}(\tau)$ and $\Real\bR_{\cI}(\tau)$,  for a simple model of the valence quark distribution
\begin{align}
\label{q-model}
  q(x) = x^{-1/2}(1-x)^3 
\end{align} 
are plotted in Fig.~\ref{fig:R-qITD}.
%
\begin{figure}[t]
\centering
\includegraphics[width=6.0cm]{BVZplots/RImodel}~
\caption{Real parts of $\cR_{\cI}(\tau)$ (thick, red) and $\bR_{\cI}(\tau)$ (thin, black) for the simple quark PDF
model in Eq.~\text{\eqref{q-model}}.
Solid curves are for $\Real R^\parallel$, dashed curves for $\Real R^\perp$.}
\label{fig:R-qITD}
\end{figure}
%
The power correction to the ``transverse'' qITD turns out
to be roughly factor four larger than for the ``longitudinal'' qITD.  
In both cases the $R$-functions flatten out at large Ioffe times $\tau \ge 10$ that are, however, hardly accessible in present day lattice
calculations. The normalization to zero proton momentum corresponds to the subtraction of the value at $\tau=0$,
so that for $z\to 0$ there is no power correction by construction. 

%
\subsection{Parton quasi-distributions}\label{sect_qPDF}%
%

Making the Fourier transformation of the above results for the qITDs we obtain the qPDFs
%
\begin{figure*}[t]
\centering
 \includegraphics[width=6.0cm]{BVZplots/RQparallel}
 ~\includegraphics[width=6.2cm]{BVZplots/RQtrans-par}
 %~\includegraphics[width=6.0cm]{BMJplots/fit-t4-qPDF}
\caption{$\cR^\parallel_\cQ(x)$ (left panel) and $\cR^\perp_\cQ(x) - \cR^\parallel_\cQ(x)$ (right panel) 
 for MSTW valence u-quarks (long black dashes), d-quarks (short blue dashes) (at 2 GeV), and for the model 
in Eq.~\eqref{q-model} (red solid curves).}
\label{fig:calR}
\end{figure*}
% 
\begin{align}
\label{t4-Qparallel}
 \cQ^\parallel(x,p) &= 
q(x) - \frac{\kappa v^2\! \Lambda^2}{(pv)^2}\Big(\frac{d}{dx} \Big)^2 \int_{|x|}^1\!\frac{dy}{y} \, (y + \bar y \ln \bar y) q(\tfrac{x}{y})
%\notag\\&= q(x) - \frac{\kappa v^2 \!\Lambda^2}{x^2(pv)^2} \int_{|x|}^1\!\frac{dy}{y} \Big[\frac{y^2}{[1-y]_+} + 2\delta(\bar y) 
%\notag\\&\qquad 
%+\delta'(\bar y)\Big]q(\tfrac{x}{y})
\notag\\&= q(x) - \frac{\kappa v^2 \!\Lambda^2}{x^2(pv)^2} 
\biggl\{
\int_{|x|}^1\!\frac{dy}{y} \frac{y^2}{[1-y]_+} q(\tfrac{x}{y})
\notag\\&\qquad
 + q(x) - |x|q'(x) \biggr\}
\end{align}
and
\begin{align}
\label{t4-Qperp}
 \cQ^\perp(x,p) &= 
q(x) - \frac{\kappa v^2\! \Lambda^2}{(pv)^2}\Big(\frac{d}{dx} \Big)^2\!\!\! \int_{|x|}^1\!\frac{dy}{y}  (y + \bar y \ln \bar y + y\bar y) q(\tfrac{x}{y})
%\notag\\&=
% q(x) - \frac{\kappa v^2 \!\Lambda^2}{x^2(pv)^2}
% \int_{|x|}^1\!\frac{dy}{y} 
%\Big[\frac{y^2}{[1-y]_+} + 3\delta(\bar y) 
%\notag\\&\qquad 
%+\delta'(\bar y)-2 y^2\Big]
%q(\tfrac{x}{y})\,.
\notag\\&= q(x) - \frac{\kappa v^2 \!\Lambda^2}{x^2(pv)^2} 
\biggl\{
\int_{|x|}^1\!\frac{dy}{y} \left[\frac{y^2}{[1-y]_+} - 2 y^2\right] q(\tfrac{x}{y})
\notag\\&\qquad
 + 2 q(x) - |x|q'(x) \biggr\}
\end{align}
Assuming for what follows $x>0$, these expressions  can be rewritten in the form
\begin{align}
\label{cR}
   \cQ^{\parallel(\perp)}(x,p) &=  q(x) \biggl\{1 - \frac{v^2 \Lambda^2}{ x^2\bar x  (pv)^2}\cR^{\parallel(\perp)}_\cQ(x) \biggr\} 
\end{align}
with
\begin{align}
 R^\parallel_Q(x) &= 
\frac{\bar x }{q(x)} \biggl\{ \int_{x}^1\!\frac{dy}{1-y} \Big[y q(\tfrac{x}{y}) - q(x)\Big] + q(x)-x q'(x)\biggr\}, 
\notag\\
 R^\perp_Q(x) &= 
\frac{\bar x }{q(x)} \biggl\{ \int_{x}^1\!\frac{dy}{1-y} \Big[y (2y-1) q(\tfrac{x}{y}) - q(x)\Big] + 2q(x)  
\notag\\&\qquad - x q'(x) \biggr\}.
\end{align}
%
Note that we have extracted the prefactor $1/(x^2\bar x)$ for the power correction anticipating that it is 
enhanced as $1/x^2$ and $1/(1-x)$ in the regions of small $x\to 0$ and large $x\to 1$ Bjorken variable, respectively.

The normalized QPDFs $\bQ^{\parallel(\perp)}(x,p)$ are obtained by replacing the kernels in the first lines
in Eqs.~\eqref{t4-Qparallel} and \eqref{t4-Qperp} by the plus distributions, cf.~\eqref{t4-nqITD}.
Writing the result in the form
\begin{align}
\label{bR}
   \bQ^{\parallel(\perp)}(x,p) &=  \hat q(x) \biggl\{1 - \frac{v^2 \Lambda^2}{ x^2\bar x  (pv)^2}\bR^{\parallel(\perp)}_\cQ(x) \biggr\}, 
\end{align}
where $\hat q(x)$ is the quark PDF normalized to the unit integral, one obtains
\begin{align}
 \bR^\parallel_\cQ(x) &= \cR^\parallel_\cQ - \frac14 \frac{\bar x x^2 }{q(x)}  q''(x)\,, %= \cR^\parallel_Q + \delta R^\parallel_\cQ\,,
\notag\\
 \bR^\perp_\cQ(x) &= \cR^\perp_\cQ - \frac5{12} \frac{\bar x x^2 }{q(x)}  q''(x)\,.    % = \cR^\perp_Q + \delta R^\perp_\cQ\,.
\end{align}
Note that the additional terms are proportional to the second derivative of the quark PDF and thus enhanced as $1/(1-x)^2$ at $x\to 1$.
As already mentioned, the normalization to the vacuum correlator does not affect the $1/p^2$ power corrections so that 
$\widehat{\bR}_{\cQ}^{\parallel(\perp)}(\tau) = \cR_{\cQ}^{\parallel(\perp)}(\tau)$ and we do not need to consider this case separately.

For a numerical study we have used the MSTW NLO valence $u$- and $d$-quark distributions \cite{Martin:2009iq} at the scale 2 GeV and the simple model
in Eq.~\eqref{q-model}. It turns out that the power corrections for (unnormalized) ``longitudinal'' and ``transverse'' qPDFs 
are similar in size so that we show $\cR^\parallel_\cQ(x)$ on the left panel in Fig.~\ref{fig:calR}, and the difference 
$\cR^\perp_\cQ(x) - \cR^\parallel_\cQ(x)$ on the right panel.

The $x$-dependence of $\cR_\cQ(x)$ is similar for all quark PDF models: The power correction is small for $x\to 0$ 
(but non-zero, $\cR^\parallel_\cQ(0)\sim 0.2$, $\cR^\perp_\cQ(0)\sim 0.4$) and 
increases steeply with $x$ (almost linearly). For the $u$-quark, the result is very similar to the simple model in Eq.~\eqref{q-model} whereas for the 
$d$-quark the power correction is, roughly, factor two larger. Note that the difference $\cR^\perp_\cQ(x) - \cR^\parallel_\cQ(x)$ depends on the 
PDF model only very weakly. 

Constructing the qPDFs from the qITDs normalized to the value at zero momentum has a large effect.
This is illustrated  in Fig.~\ref{fig:bRparallel} where in the upper panels we show the results for the ``longitudinal'' case and in 
the lower panels the difference between the ``longitudinal'' and ``transverse'' distributions.  
%
\begin{figure*}[tbp]
\centering
 \includegraphics[width=5.0cm]{BVZplots/normRQu}~
 \includegraphics[width=5.0cm]{BVZplots/normRQd}~
 \includegraphics[width=5.0cm]{BVZplots/normRQm}
\\
 \includegraphics[width=5.0cm]{BVZplots/deltaR-RQu}~
 \includegraphics[width=5.0cm]{BVZplots/deltaR-RQd}~
 \includegraphics[width=5.0cm]{BVZplots/deltaR-RQm}
\caption{Upper panels:
$\cR^\parallel_\cQ$ (solid red curves)  and $\bR^\parallel_\cQ$ (dashed black curves) 
for MSTW valence u-quarks (left), d-quarks (middle), both at 2 GeV, and for the  model 
in Eq.~\eqref{q-model} (right).
Lower panels: the same for 
$\cR^{\perp}_\cQ-\cR^{\parallel}_\cQ $ and $\bR^{\perp}_\cQ-\bR^{\parallel}_\cQ $.
Note split panels: the results for $0<x<0.6$ and $0.6<x<1$ are shown using a different scale on the vertical axis.}
\label{fig:bRparallel}
\end{figure*}
%
The three panels (from left to right) correspond to the MSTW valence u-quarks, d-quarks, and the simple model 
in Eq.~\eqref{q-model}, respectively. The red solid lines stand for  $\cR_\cQ(x)$ 
(i.e. the same as in Fig.~\ref{fig:calR}), and the result for the normalized qPDFs, $\bR_\cQ(x)$, is shown by the black dashed curves.

We see that the normalization to the qITD at zero momentum significantly reduces the power correction at moderate values of $x \lesssim 0.5$
at the cost of dramatical increase at higher $x$ values. This normalization procedure is, therefore, not suitable to access the large-$x$ behavior
of the PDFs, but apparently minimizes power corrections in the not-so-large $x$ region. 

%
\subsection{Parton pseudo-distributions}\label{sect_pPDF}%
%

\begin{figure}[t]
\centering
 \includegraphics[width=5.0cm]{BVZplots/pseudoR}~
% \includegraphics[width=5.0cm]{BMJplots/t4-pPDF}~
\caption{\label{fig:pPDF} Power correction to the pPDF $\mathcal{R}_\cP$ \eqref{t4-pPDF}
for MSTW valence  u-quarks (long black dashes), d-quarks (short blue dashes), both at 2 GeV, and for the simple model 
$q(x) = x^{-1/2} (1-x)^3 $ (red solid curve).}
\end{figure}

Power corrections for the pPDF can be obtained easily from the corresponding expression for the 
``transverse'' qITDs \eqref{t4-qITD}, \eqref{t4-nqITD}. 
Writing the result as
\begin{align}
\mathcal{P}(x,z,\mu) &= q(x)\Big\{1 + (v^2\! z^2\!\Lambda^2)  \theta(|x|<1) \mathcal{R}_\cP(x)\Big\},
\end{align}
and similarly for the pPDF normalized to zero momentum, $\mathbf{P}(x,z)$, we obtain
\begin{align}
\label{t4-pPDF}
\mathcal{R}_\cP(x) &= \frac{1}{q(x)} \int_{|x|}^1\!\frac{dy}{y} \, (y + \bar y \ln \bar y +  y \bar y ) q(\tfrac{x}{y})\,,
\notag\\
\mathbf{R}_\cP(x) &= \frac{1}{q(x)} \int_{|x|}^1\!\frac{dy}{y} \, [y + \bar y \ln \bar y +  y \bar y ]_+ q(\tfrac{x}{y}) 
\notag\\&= 
\mathcal{R}_\cP(x) - 5/12 \,.
\end{align}
The numerical results are shown in Fig.~\ref{fig:pPDF}.

Note that $\mathcal{R}_\cP(x)$ is very similar for all considered models for the valence quark PDFs and
decreases at $x\to 1$. Indeed, it is easy to see that $\mathcal{R}_\cP(x) = \mathcal{O}(1-x)$ in this limit, 
similar to the target mass correction, Eq.~\eqref{target-pPDF}. This suppression is removed once pPDF is
normalized to zero momentum ~\cite{Orginos:2017kos} (which adds a negative constant), but can be upheld if
normalized to the vacuum expectation value of the  same operator. 

%
